\subsection{Estados Financieros}


\textbf{Compañía de Seguros: Quálitas}

En el caso de Quálitas, el riesgo de mercado se entiende como la pérdida potencial derivada de cambios en factores que afectan el valor de sus activos y pasivos, tales como tasas de interés, tipo de cambio e índices de precios. La propia compañía lo define así dentro de su sistema de administración integral de riesgos. Esta exposición resulta especialmente relevante porque una parte central de su balance está concentrada en inversiones en valores, las cuales al cierre de 2024 ascendieron a \$47{,}821.5 millones, principalmente en instrumentos gubernamentales, deuda privada, renta variable y valores extranjeros.

La importancia de este riesgo también se observa en el estado de resultados. Durante 2024, Quálitas registró dentro del resultado integral de financiamiento ingresos por inversiones por \$3{,}292.7 millones, pérdidas por venta de inversiones de \$167.7 millones, ganancias por valuación de inversiones de \$1{,}030.7 millones y un resultado cambiario positivo de \$101.4 millones. Esto muestra que una fracción importante de su utilidad depende del desempeño de los mercados financieros y no solamente de la operación técnica del seguro. En otras palabras, si hubiera movimientos adversos en tasas o en el tipo de cambio, el valor de su portafolio podría deteriorarse y afectar directamente la utilidad neta.

Respecto al riesgo de crédito, Quálitas lo define como la pérdida potencial derivada de la falta de pago o del deterioro en la solvencia de contrapartes y deudores. La empresa señala que este riesgo se origina principalmente en las cuentas por cobrar y en las inversiones en instrumentos de deuda, además de incluir exposiciones asociadas a contratos de reaseguro y cuentas por cobrar de intermediarios. En 2024, uno de los rubros más relevantes del activo fue el deudor por primas, con un saldo de \$40{,}089.8 millones, además de una cartera de crédito vigente de \$765.9 millones, cartera vencida por \$30.5 millones y estimaciones preventivas para riesgo crediticio por \$139.3 millones. Estos datos reflejan que una parte considerable de los activos de la aseguradora depende de la recuperación efectiva de cobros a clientes, agentes y otras contrapartes.

Adicionalmente, Quálitas reconoce deterioros vinculados al riesgo de crédito en resultados. Por ejemplo, en 2024 registró un incremento por deterioro en instrumentos de deuda corporativa con cambios en la utilidad por \$7.5 millones, y también aplica matrices de pérdida crediticia esperada para otras cuentas por cobrar. Esto muestra que la entidad no solo reconoce el riesgo en forma conceptual, sino que también lo mide y lo refleja contablemente mediante provisiones y castigos preventivos. Así, el riesgo de crédito en Quálitas no se limita a la cartera tradicional, sino que se extiende a toda la calidad financiera de sus contrapartes.

En cuanto al riesgo de liquidez, Quálitas lo define como la pérdida potencial derivada de la venta anticipada o forzosa de activos con descuentos inusuales para hacer frente a sus obligaciones, o por la imposibilidad de enajenar oportunamente una posición. En una aseguradora este riesgo es especialmente importante, porque la compañía debe mantener recursos suficientes para hacer frente al pago de siniestros, reservas técnicas y otras obligaciones operativas. En el balance de 2024 se observa que Quálitas contaba con efectivo y equivalentes por \$3{,}952.7 millones, pero al mismo tiempo enfrentaba pasivos significativos, entre ellos reservas técnicas por más de \$59{,}301 millones y acreedores por \$11{,}432.3 millones. Esto implica que la administración de la liquidez depende en gran medida de la calidad y disponibilidad de sus inversiones, así como del adecuado calce entre activos y obligaciones.

La relevancia de la liquidez también se aprecia en el marco regulatorio que enfrenta la institución. Quálitas señala que está sujeta a requerimientos de Base de Inversión y Requerimiento de Capital de Solvencia (RCS), cuyo objetivo es asegurar que cuente con recursos patrimoniales suficientes para hacer frente tanto a sus riesgos asumidos como a condiciones excepcionales de mercado. Además, las reservas técnicas se constituyen precisamente para cubrir obligaciones futuras derivadas del pago de siniestros y otros compromisos contractuales. Esto significa que la liquidez en una aseguradora no puede analizarse solo desde el efectivo disponible, sino desde la capacidad estructural de convertir activos en recursos oportunos sin comprometer la estabilidad financiera de la empresa.

En conjunto, el caso de Quálitas muestra una estructura de riesgos típica de una aseguradora grande: el riesgo de mercado se concentra en la valuación de su portafolio de inversiones y en los efectos de variables financieras sobre el resultado integral de financiamiento; el riesgo de crédito se encuentra en la exposición a deudores por primas, cartera de crédito, reaseguradores y otros deudores; y el riesgo de liquidez se relaciona con la necesidad de respaldar reservas técnicas y cumplir oportunamente con siniestros y demás obligaciones. Por ello, el análisis de sus estados financieros permite ver que la rentabilidad y la solvencia de la compañía dependen no solo de vender pólizas, sino de administrar adecuadamente estos tres riesgos financieros.

\newpage
\textbf{Grupo Financiero Banorte}

En el caso de Grupo Financiero Banorte, el riesgo de cr\'edito es uno de los m\'as relevantes, debido a que una parte muy importante de su operaci\'on se basa en la colocaci\'on de financiamiento y en la administraci\'on de activos sujetos a incumplimiento. Este riesgo puede entenderse como la posibilidad de p\'erdida derivada de que los acreditados o contrapartes no cumplan con sus obligaciones de pago. En los estados financieros de 2024, este riesgo se refleja con claridad en la cartera de cr\'edito total, que ascendi\'o a \$1{,}190{,}781 millones antes de reservas y a \$1{,}178{,}167 millones netos despu\'es de la estimaci\'on preventiva. Asimismo, Banorte muestra una cartera clasificada en etapa 1, etapa 2 y etapa 3, lo cual permite distinguir entre cr\'editos con bajo riesgo, con deterioro significativo y con incumplimiento m\'as evidente. A ello se suma una estimaci\'on preventiva para riesgos crediticios de \$20{,}122 millones, lo que confirma la relevancia de este riesgo dentro de su balance.

La importancia del riesgo de cr\'edito tambi\'en se observa en el estado de resultados. Durante 2024, Banorte report\'o un margen financiero de \$139{,}360 millones, pero despu\'es de reconocer la estimaci\'on preventiva para riesgos crediticios por \$19{,}536 millones, el margen financiero ajustado se redujo a \$119{,}824 millones. Esto implica que una parte significativa del resultado depende de la calidad de la cartera y de la probabilidad de que los clientes paguen oportunamente. Adem\'as, el propio documento indica que la medici\'on de este riesgo se basa en variables como la probabilidad de incumplimiento, la severidad de la p\'erdida y la exposici\'on al incumplimiento, las cuales son esenciales para determinar las reservas que Banorte debe constituir.

Por su parte, el riesgo de mercado en Banorte se relaciona con la posibilidad de p\'erdidas ocasionadas por movimientos adversos en tasas de inter\'es, tipo de cambio y precios de instrumentos financieros. Este riesgo es particularmente importante porque el grupo mantiene una posici\'on considerable en inversiones en instrumentos financieros, las cuales al cierre de 2024 sumaron \$976{,}597 millones. Dentro de este rubro se incluyen instrumentos negociables, instrumentos para cobrar o vender y valores para cobrar principal e inter\'es. Asimismo, Banorte registr\'o activos por derivados por \$22{,}130 millones, lo que muestra que participa activamente en operaciones financieras sensibles a cambios de mercado.

La exposici\'on al riesgo de mercado tambi\'en se hace visible en el estado de resultados y en el resultado integral. En 2024, Banorte obtuvo un resultado por intermediaci\'on de \$4{,}983 millones, adem\'as de registrar efectos por valuaci\'on de instrumentos financieros para cobrar o vender y por derivados de cobertura. En el otro resultado integral aparecen p\'erdidas por valuaci\'on de instrumentos financieros para cobrar o vender de \$1{,}783 millones y por valuaci\'on de instrumentos financieros derivados de cobertura de flujos de efectivo de \$1{,}047 millones. Esto muestra que la rentabilidad del grupo no depende \'unicamente de la actividad crediticia tradicional, sino tambi\'en del comportamiento de los mercados financieros y de la valuaci\'on de sus posiciones de tesorer\'ia e inversi\'on.

En cuanto al riesgo de liquidez, Banorte enfrenta la necesidad de cumplir oportunamente con sus obligaciones financieras sin tener que vender activos en condiciones desfavorables ni recurrir a fondeo excesivamente costoso. En una instituci\'on bancaria este riesgo es especialmente relevante, porque una gran parte de su estructura de pasivos se sustenta en dep\'ositos del p\'ublico, emisiones y otras fuentes de financiamiento que deben ser administradas con precisi\'on. En el balance de 2024 se observa que Banorte ten\'ia efectivo y equivalentes de efectivo por \$98{,}704 millones, pero al mismo tiempo manten\'ia una captaci\'on tradicional de \$1{,}133{,}500 millones, integrada principalmente por dep\'ositos de exigibilidad inmediata y dep\'ositos a plazo. Esta diferencia muestra que la liquidez del banco no puede medirse solo con el efectivo disponible, sino con su capacidad de transformar activos financieros en recursos l\'iquidos y con la estabilidad de su fondeo.

Las notas de los estados financieros refuerzan esta idea al se\~nalar que, en moneda extranjera, Banorte estuvo sujeto al r\'egimen de liquidez de Banco de M\'exico. Durante 2024, la instituci\'on gener\'o un requerimiento promedio mensual de liquidez de 964{,}178 miles de USD, mientras que mantuvo una inversi\'on promedio en activos l\'iquidos de 5{,}324{,}763 miles de USD, con un exceso promedio de 4{,}360{,}584 miles de USD. Esto sugiere que Banorte conserv\'o una posici\'on holgada para atender sus obligaciones en moneda extranjera. Al mismo tiempo, la composici\'on de la captaci\'on tradicional muestra que el banco depende ampliamente de dep\'ositos a la vista y a plazo, por lo que la adecuada administraci\'on de vencimientos y flujos de efectivo resulta fundamental para sostener su operaci\'on.

En conjunto, el caso de Banorte muestra una estructura de riesgos t\'ipica de una instituci\'on bancaria de gran tama\~no. El riesgo de cr\'edito domina por el volumen de su cartera y por la necesidad de constituir reservas preventivas que impactan directamente el resultado. El riesgo de mercado se concentra en la valuaci\'on de inversiones, derivados y dem\'as posiciones financieras expuestas a cambios en tasas y precios. Finalmente, el riesgo de liquidez se relaciona con la administraci\'on del efectivo, de la captaci\'on tradicional y de los activos l\'iquidos necesarios para cubrir obligaciones inmediatas y regulatorias. Por ello, el an\'alisis de sus estados financieros permite concluir que la solidez de Banorte depende no solo de su capacidad para generar utilidades, sino tambi\'en de la administraci\'on integral de estos tres riesgos financieros.

\newpage
\textbf{Monex Casa de Bolsa}

En el caso de Monex Casa de Bolsa, el riesgo de cr\'edito se entiende como la posibilidad de p\'erdida derivada del incumplimiento de una contraparte o acreditado en las operaciones que realiza la instituci\'on. De acuerdo con sus notas, este riesgo se concentra principalmente en las operaciones de reporto y en los derivados extraburs\'atiles con clientes. La propia casa de bolsa se\~nala que, para mitigar esta exposici\'on, establece l\'ineas de operaci\'on con base en el an\'alisis financiero de cada contraparte y exige garant\'ias que van de 6\% a 10\% del monto de la operaci\'on, adem\'as de aplicar llamadas de margen cuando las posiciones abiertas presentan p\'erdidas relevantes. Tambi\'en aclara que su portafolio de deuda no presenta pr\'acticamente riesgo de cr\'edito emisor, ya que opera fundamentalmente con t\'itulos gubernamentales.

Este riesgo puede observarse en el balance a trav\'es de rubros como deudores por reporto por \$48{,}226 millones, otras cuentas por cobrar por \$37{,}617 millones e instrumentos financieros derivados con saldo activo por \$1{,}235 millones. Estas cuentas muestran que una parte importante del activo depende de la capacidad de pago y del cumplimiento de contrapartes financieras y clientes. Por ello, aunque Monex opera con mecanismos de mitigaci\'on, sigue existiendo exposici\'on crediticia en sus actividades de intermediaci\'on y derivados.

Por su parte, el riesgo de mercado en Monex se refiere a la p\'erdida potencial originada por cambios adversos en factores como tasas de inter\'es, tipos de cambio e \'indices de precios, que afectan la valuaci\'on de sus posiciones y los resultados esperados. La instituci\'on eval\'ua este riesgo mediante modelos de Valor en Riesgo (VaR) con un horizonte de 1 d\'ia y un nivel de confianza de 99\%, complementados con pruebas de estr\'es. Al cierre de 2024, el VaR global fue de \$7.79 millones, equivalente a 1.01\% del capital neto, lo que permite dimensionar la p\'erdida m\'axima esperada bajo condiciones normales de mercado en el corto plazo.

La exposici\'on al riesgo de mercado se refleja claramente en el balance y en el estado de resultados. En el activo, Monex report\'o inversiones en instrumentos financieros por \$1{,}101 millones y derivados por \$1{,}235 millones, ambos sensibles a variaciones de mercado. En el estado de resultados, el resultado por intermediaci\'on fue de \$287 millones en 2024, integrado por partidas muy ligadas a la volatilidad financiera, como una ganancia por valuaci\'on de derivados de \$1{,}172 millones, una p\'erdida por valuaci\'on de divisas de \$1{,}150 millones, una utilidad por compraventa de instrumentos negociables de \$171 millones y otra por compraventa de derivados de \$60 millones. Esto muestra que la rentabilidad de la casa de bolsa depende en gran medida del comportamiento de los mercados financieros, especialmente en negocios de divisas, derivados y productos burs\'atiles.

En cuanto al riesgo de liquidez, Monex lo define como la posibilidad de no poder atender oportunamente sus necesidades de flujo de efectivo, no renovar pasivos en condiciones normales o verse obligada a vender activos anticipadamente con descuentos inusuales. La instituci\'on tambi\'en reconoce que este riesgo puede surgir por descalces de plazo entre activos y pasivos. Para su administraci\'on, monitorea continuamente la liquidez de los portafolios de mercado de dinero y derivados, as\'i como el perfil de vencimientos de las operaciones de reporto.

Este riesgo se aprecia con claridad en la estructura del balance. Al cierre de 2024, Monex ten\'ia efectivo y equivalentes por solo \$183 millones, frente a pasivos y compromisos relevantes como colaterales vendidos o dados en garant\'ia por \$48{,}072 millones, acreedores diversos y otras cuentas por pagar por \$37{,}917 millones, y acreedores por liquidaci\'on de operaciones por \$932 millones. Adem\'as, el cuadro de vencimientos muestra que gran parte de los activos y pasivos est\'a concentrada en plazos de hasta 6 meses, lo que obliga a una administraci\'on muy cuidadosa de flujos de corto plazo. Aun as\'i, la propia instituci\'on reporta un indicador de liquidez de 225.66\% en 2024, lo que sugiere una posici\'on relativamente holgada para enfrentar sus obligaciones l\'iquidas inmediatas.

Adicionalmente, las operaciones de reporto son un elemento central en la liquidez de Monex. La casa de bolsa report\'o deudores por reporto por \$48{,}226 millones y acreedores por reporto por \$274 millones, con plazos promedio relativamente cortos. Los intereses a favor por estas operaciones ascendieron a \$6{,}521 millones, mientras que los intereses a cargo fueron de \$6{,}344 millones en 2024. Esto revela que una parte importante de su operaci\'on diaria depende del fondeo y renovaci\'on eficiente de reportos, por lo que cualquier tensi\'on en este mercado podr\'ia afectar directamente su liquidez y su margen financiero.

En conjunto, el caso de Monex Casa de Bolsa muestra una estructura de riesgos propia de una instituci\'on de intermediaci\'on financiera enfocada en mercado de dinero, divisas, derivados y productos burs\'atiles. El riesgo de cr\'edito se concentra en contrapartes de reportos y derivados OTC; el riesgo de mercado surge de la valuaci\'on de inversiones, derivados y posiciones cambiarias, y se mide con VaR y pruebas de estr\'es; mientras que el riesgo de liquidez depende de la administraci\'on de flujos de corto plazo, reportos, colaterales y otras cuentas por pagar. Por ello, el an\'alisis de sus estados financieros permite concluir que la estabilidad y rentabilidad de Monex dependen no solo del volumen de sus operaciones, sino de una administraci\'on estricta e integral de estos tres riesgos